\capitulo{2}{Objetivos del proyecto}

Este apartado explica de forma precisa y concisa cuales son los objetivos que se persiguen con la realización del proyecto. Se puede distinguir entre los objetivos marcados por los requisitos del software a construir y los objetivos de carácter técnico que plantea a la hora de llevar a la práctica el proyecto.

Los Grandes Modelos de Lenguaje (en inglés, Large Language Models, LLM) constituyen una de las principales innovaciones que recientemente han eclosionado en el campo del procesamiento del lenguaje natural. Estos modelos, basados en arquitecturas de transformers y entrenados con grandes volúmenes de datos textuales, son capaces de capturar relaciones sintácticas y semánticas complejas, lo que les permite abordar tareas de generación, resumen, traducción y razonamiento en lenguaje natural. No obstante, presentan limitaciones inherentes: el conocimiento que contienen está restringido a los datos utilizados en su entrenamiento y a una fecha de corte determinada, lo que puede derivar en respuestas desactualizadas, incompletas o incluso incorrectas en dominios especializados.

En este contexto, la técnica denominada Retrieval-Augmented Generation (RAG) surge como una estrategia para superar dichas limitaciones. Los RAG combinan un modelo generativo con un sistema externo de recuperación de información (de ahí proviene su nombre), de modo que la generación textual no se sustenta únicamente en el conocimiento latente del LLM, sino que se apoya en la consulta de un repositorio documental (bases de datos vectoriales). A través de este enfoque híbrido, el modelo es capaz de producir respuestas más precisas y fundamentadas, al incorporar fragmentos de información verificados y actualizados provenientes de bases de datos o colecciones textuales específicas.

Durante la beca de colaboración se pretende que el estudiante entienda y llegue a implementar un RAG, este proceso comprende varias etapas. Inicialmente, es necesario realizar un preprocesamiento de la colección documental, lo que incluye tareas de segmentación en fragmentos, normalización del texto y generación de representaciones vectoriales mediante técnicas de embeddings. Posteriormente, estos vectores se almacenan en un índice semántico o base de datos vectorial (como FAISS, Milvus o Pinecone), optimizada para consultas de similitud. Esta infraestructura constituye el núcleo del mecanismo de recuperación de contexto.

Finalmente, en la fase de interacción con el usuario, la consulta planteada se transforma también en una representación vectorial, la cual se utiliza para identificar los fragmentos más relevantes dentro del índice. Estos fragmentos se integran en el prompt proporcionado al LLM, de manera que el modelo genera una respuesta enriquecida y contextualizada. Este esquema dota a los LLM de la capacidad de adaptarse dinámicamente a distintos dominios de conocimiento, incrementando su fiabilidad y ampliando su aplicabilidad en entornos de investigación e innovación tecnológica.

